\chapter{Introduction}
\label{chap:introduction}

\section{The measurement problem}

Reading out a superconducting qubit means detecting a microwave signal of a few
photons against the thermal and amplifier noise of the receiving chain. The
first amplifier in that chain sets the noise floor for everything downstream:
if it adds noise, no later stage can remove it. A high-electron-mobility
transistor amplifier, the conventional first stage, adds of order ten quanta of
noise. That is enough to force long integration times and, for multiplexed
readout of many qubits, to make single-shot fidelity unreachable.

A parametric amplifier built from Josephson junctions can in principle reach the
quantum limit: half a photon of added noise for a phase-insensitive amplifier,
and formally zero for a phase-sensitive one. The mechanism is not transistor
transconductance but \emph{parametric conversion}: a strong pump modulates a
reactance, and the modulated reactance transfers energy from the pump into a
weak signal.

The difficulty is bandwidth. A resonant Josephson parametric amplifier (JPA)
concentrates the pump in a cavity and buys large gain over a bandwidth of a few
tens of megahertz --- enough for one qubit, not for twenty. The
\emph{travelling-wave} parametric amplifier (TWPA) removes the cavity: the
signal propagates down a long nonlinear transmission line, accumulating gain
continuously. Bandwidths of several gigahertz become possible. The price is that
the device is now a distributed system with thousands of degrees of freedom, and
its behaviour is governed by phase matching along the line rather than by a
single resonance condition.

\section{Why this is a hard simulation problem}

A TWPA is a chain of $10^{3}$--$10^{4}$ unit cells, each containing at least one
Josephson junction. Three properties conspire to make it awkward:

\begin{enumerate}
\item \textbf{It is strongly driven.} The pump is not a perturbation. At the
      operating point the junction phase excursion is a substantial fraction of
      a radian, so the pump waveform contains significant harmonic content and
      the small-signal expansion of the junction nonlinearity is not adequate.
      The pump must be found as a genuine \emph{periodic steady state} of a
      nonlinear system, not as a linear response.

\item \textbf{It is stiff and distributed.} The chain couples nearest
      neighbours only, which is good --- the matrices are banded --- but the
      condition number grows with length, and the physically interesting
      operating points sit close to the boundary where the periodic solution
      ceases to exist. Naive Newton iteration from a cold start diverges.

\item \textbf{The interesting observables live at different orders.} Gain is a
      \emph{linear} response about the pumped state. Compression is not: at
      large signal power the generated sidebands react back on the pump, and
      pump and signal must be solved simultaneously. The two calculations have
      radically different cost, and conflating them is the single most common
      way to produce a saturation number that is really a small-signal number
      in disguise.
\end{enumerate}

Time-domain integration is possible but poorly suited: one must integrate for
many pump periods to reach steady state, then again for every signal frequency,
and the transient decay time of a low-loss superconducting line is long. The
natural method is \emph{harmonic balance}: assume periodicity, expand in a
finite Fourier basis, and solve the resulting nonlinear algebraic system
directly. This thesis is about doing that correctly, and about what the result
does and does not establish.

\section{Scope and contributions}

The work described here is the design, implementation and validation of
\code{twpa\_solver}, a harmonic-balance solver for Josephson circuits, together
with the physical study it enables. The specific contributions are:

\begin{enumerate}
\item \textbf{A three-stage solver architecture} (\cref{chap:pumphb},
      \cref{chap:floquet}, \cref{chap:saturation}) that separates the nonlinear
      pump solve, the linear small-signal Floquet solve, and the fully coupled
      multitone solve. The separation is what makes a $50\times50$ gain map
      affordable while keeping a route to saturation that does not assume the
      signal is small.

\item \textbf{An exact real-packed preconditioner} for the harmonic-balance
      Jacobian (\cref{sec:real-coupled}) that retains the conjugate coupling
      term dropped by the conventional complex mode-coupled construction.
      Because the residual is not complex-analytic in the phasor variables,
      dropping that term costs accuracy exactly where the device is stiffest;
      retaining it makes GMRES converge in about one iteration.

\item \textbf{A continuation and recovery layer} (\cref{sec:continuation})
      built around natural-parameter continuation in the source amplitude, with
      secant and tangent predictors, adaptive step control, and a resume-rather-
      than-restart fallback.

\item \textbf{A two-dimensional tone-lattice formulation of saturation}
      (\cref{chap:saturation}) in which pump harmonics and signal sidebands are
      indexed by a single integer pair $(h,q)$ on a torus, with exact FFT-based
      transforms and a basis constructor that provably covers a requested
      Floquet sideband range.

\item \textbf{A careful account of what is and is not validated}
      (\cref{chap:validation}). Several results in this field are validated
      against another simulator; this thesis argues that such agreement is a
      regression check and not physical evidence, and states plainly which
      quantities currently rest on no external reference.
\end{enumerate}

\section{Structure of this thesis}

\Cref{chap:theory} develops the physics: the Josephson element, the
travelling-wave geometry, three- and four-wave mixing, phase matching, and the
coupled-mode picture that gives the standard analytic gain expression. It ends
by identifying precisely where that picture breaks down and why a numerical
method is required.

\Cref{chap:circuit} sets up the circuit model: the node-flux formulation, the
matrix stamps, the treatment of the nonlinear branches, and the design families
the solver handles.

\Cref{chap:pumphb} is the core algorithmic chapter. It develops the harmonic-
balance residual from first principles and specifies every nested loop of the
solve --- frequency, power, continuation, Newton, residual evaluation,
preconditioner assembly, factorisation, Krylov solve, Jacobian-vector product,
and line search --- at the level of detail required to reimplement it. Where
the implementation departs from the idealised statement, the departure is
recorded explicitly in an \emph{Implementation} box giving the actual coded
form.

\Cref{chap:numerics} treats the numerical linear algebra separately: Schur
reduction, the preconditioner hierarchy, factorisation backends, ordering, and
the measured performance consequences of each.

\Cref{chap:floquet} develops the small-signal regime: the Floquet picture, the
sideband conversion matrix, gain definitions, and quantum efficiency.

\Cref{chap:saturation} develops the saturation regime: the tone lattice, the
multitone residual, compression curves, $P_{1\mathrm{dB}}$, power balance,
Manley--Rowe invariants, and the stability of the large-signal state.

\Cref{chap:validation} states the validation position honestly, including the
retraction of a previously used reference.

\Cref{chap:conventions} is a reference appendix of conventions and the specific
errors each one guards against.
